\documentclass[a4paper, eqnumbering=section, thmnumbering=section, 10pt]{simple-eng}

% \renewcommand{\cftsecleader}{\hfill} % tocにおいて，セクション名の右のドットを排除する


\title{September 10, 2025}
\author{Yuki Yamamoto}
\date{\today}



\begin{document}
\maketitle
\tableofcontents

\newpage
\begin{myproof}
    its okasiii!!
\end{myproof}
\begin{myproof}
    its okassiiiii\\
    lkakfdaofjpi
    \begin{align}
        \dfrac{1}{\int_{\alpha}}
    \end{align}
    $\dfrac{1}{\int_{\alpha}}$
\end{myproof}
\begin{definition}{}{}
    its okasiiii!!!
\end{definition}
\begin{definition}
    its okassiiiii\\
    lkakfdaofjpi
    \begin{align}
        \dfrac{1}{\int_{\alpha}}
    \end{align}
    $\dfrac{1}{\int_{\alpha}}$
\end{definition}

\newpage
Let us fix an algebraically closed field $k$ of arbitrary charactristic.

\setcounter{section}{-1}
\section{Notations and Preliminaries}



\section{Self-Intersection Number}

\begin{definition}{}{}
    Let $X$ be a nonsingular variety over $k$, $Y$ be a nonsingular subvaraiety of $X$, 
    $\textrm{codim}(Y;X)=r$, and for an ideal shaef $\mathscr{I}$ of $\mathscr{O}_X$, 
    $Y=Z(\mathscr{I})$. 
    \begin{enumerate}[label=(\arabic*{})]
        \item $\mathscr{I}/\mathscr{I}^2$ is a locally free $\mathscr{O}_Y$-module of rank $r$, 
            and is called \textit{\textbf{conormal sheaf}} of $Y$.
        \item $\mathscr{N}_{Y/X} \coloneqq (\mathscr{I}/\mathscr{I}^{2})^{\vee }$ is 
            called \textbf{\textit{normal sheaf} }of $Y$.
    \end{enumerate}
\end{definition}


\begin{example}{}{}
    Let $X$ be a nonsingular variety, and $D$ be any divisor on $X$. The 
    \textbf{\textit{self-intersection number}}
    $D^2=D.D$ can be calculated as follows.\par
    \begin{itemize}
        \item If $D$ is very ample, then by (1.2), we can take an irreducible nonsingular curve 
            $C,C^\prime \in \left\lvert D \right\rvert $ which meet transvarsally each other, and
            \begin{align}
                D.D = C.C^\prime &=\text{deg}_C(\mathscr{O}_X(C^\prime)\otimes \mathscr{O}_C)\\
                &= \text{deg}_C(\mathscr{O}_X(C)\otimes \mathscr{O}_C)
            \end{align}
        \item $(\mathscr{O}_X(C)\otimes \mathscr{O}_C)^{\vee } = 
            \mathscr{O}_X(-C)\otimes \mathscr{O}_C = (\mathscr{N}_{C/X})^{\vee}$, thus $D^2 = \text{deg}_C(\mathscr{N}_{C/X})$
        \item For any divisor $D$, we can choose nonsingular and very ample curve $C,C^\prime$ such that 
            $D \sim C-C^\prime $
    \end{itemize}
\end{example}

\section{Riemann-Roch Theorem for a Surface}

\begin{definition}{}{}
    Let $X$ be a nonsingular variety$/k$, and $n=\dim X$.
    \begin{enumerate}[label=(\arabic*{})]
        \item $p_g \coloneqq \dim_k H^0(X,\omega_X) = h^0(\omega_X)$ is called the 
            \textbf{\textit{geometric genus}} of $X$.
        \item $p_a = (-1)^n (\chi(\mathscr{O}_X) - 1)$ is called the \textbf{\textit{arithmetic genus}} of $X$.
    \end{enumerate}
    If $n=1$, by Serre duality, 
    \begin{align}
        p_a = 1-\chi(\mathscr{O}_X) = h^1(\mathscr{O}_X) = h^0(\omega_X) = p_g \; . 
    \end{align}
    If $n=2$, 
    \begin{align}
        p_a &= h^2(\mathscr{O}_X)- h^1(\mathscr{O}_X) = h^0(\omega_X) - h^1(\omega_X) \\
        &= p_g - h^1(\omega_X)\;. 
    \end{align}
\end{definition}

Here we briefly recall the Riemann–Roch theorem for curves.
Let $C$ be a cruve of genus $g$, and $D\in \text{div}(C)$. The formula is given by
\begin{align}\label{eq:riemann-roch for a curve}
    \chi(\mathscr{O}_C(D)) = \deg D + 1 - g
\end{align}
By Serre duality, the left hand side is also expressed by 
$\chi(\mathscr{O}_C(D)) = h^0(\mathscr{O}_C(D)) - h^0(\mathscr{O}_C(K-D)) $, where $K$ is 
the canonical divisor on $X$. In perticular, if $D=K$, \eqref{eq:riemann-roch for a curve} 
is rewriten as 
\begin{align}\label{eq:if D=K(curve)}
    2g-2 = \deg_C(\omega_C) 
\end{align}

\begin{proposition}{Adjunction Formula}{Adjunction Formula}
Let $X$ be a surface, and $K$ be the canonical divisor on $X$.
For any nonsingular curve $C$ of genus $g$ on $X$, we have;
\begin{align}
    2g-2 = C.(C+K) \;.
\end{align}
\end{proposition}
\begin{myproof}
    Let $\omega_X, \omega_C$ be the canonical sheaves on $X, C$ each other. 
    From [II. 8.20], we have 
    \begin{align}
    \omega_C 
    &\cong \left(\omega_X\otimes_{\mathscr{O}_X}\mathscr{O}_X(C)\right)\otimes\mathscr{O}_C\\
    &\cong \mathscr{O}_X(C+K)\otimes \mathscr{O}_C
    \end{align}
    Then, by \eqref{eq:if D=K(curve)}, 
    \begin{align}
        2g-2 &= \deg_C(\omega_C) = \deg_C(\mathscr{O}_X(C+K)\otimes \mathscr{O}_C)\\
            &= C.(C+K)
    \end{align}\hfill
\end{myproof}

\begin{remark}
    For any Noetherian and separated schemes $X,Y$, an affine morphism $f:X\ri Y$, and 
    a quasi-coherent $\mathscr{O}_X$-module $\mathscr{F}$, we have 
    $H^i(X,\mathscr{F}) \cong H^i(Y,f_\ast\mathscr{F})$. 
\end{remark}

\begin{theorem}{Riemann-Roch Theorem}{}
    Let $X$ be a surface, and $K$ be the canonical divisor on $X$. 
    Then for any divisor $D$ on $X$, we have
    \begin{align}
        \chi(\mathscr{O}_X(D)) = \dfrac{1}{2}D.(D-K) + p_a + 1
    \end{align}
\end{theorem}
\begin{myproof}
    Let $D\sim C-E$, where $C,E$ are smooth curves. Then, we have two exact sequences of 
    $\mathscr{O}_X$-modules:
    \[0 \ri \mathscr{O}_X(-C)\ri \mathscr{O}_X \ri \mathscr{O}_X/\mathscr{O}_X(-C) \ri 0, \] and
    \[ 0 \ri \mathscr{O}_X(-E)\ri \mathscr{O}_X \ri \mathscr{O}_X/\mathscr{O}_X(-E) \ri 0. \]
    Tensoring with $\calO_X(C)$, we obtain
    \[
        0 \ri \mathscr{O}_X \ri \mathscr{O}_X(C) \ri \mathscr{O}_X\otimes (\mathscr{O}_X/\mathscr{O}_X(-C)) \ri 0,
    \]
    and
    \[
        0 \ri \mathscr{O}_X(C-E) \ri \mathscr{O}_X(C) \ri \mathscr{O}_X(C)\otimes (\mathscr{O}_X/\mathscr{O}_X(-E))\ri 0.
    \]
    Let $f:C \hookrightarrow X$ and $g:E\hookrightarrow X$ denote the closed immersions. By 
    additivity of the Euler-Poincaré characteristic, 
    \setlength{\jot}{10pt} 
    \begin{align}
        \chi(\calO_X(D)) &= \chi(\mathscr{O}_X(C-E)) \\ 
        &= \chi(\mathscr{O}_X(C)) - 
        \chi(\mathscr{O}_X(C)\otimes(\mathscr{O}_X/\mathscr{O}_X(-E)))\\
        &=  \chi\left(\mathscr{O}_X(C)\otimes (\mathscr{O}_X/\mathscr{O}_X(-C))\right)
        - \chi\left( \mathscr{O}_X(C)\otimes(\mathscr{O}_X/\mathscr{O}_X(-E)) \right)
        + \chi(\calO_X)\\
        &= \chi(f_\ast f^\ast \calO_X(C)) - \chi(g_\ast g^\ast \calO_X(C)) + p_a + 1\\
        &= \chi(f^\ast \calO_X(C)) - \chi(g^\ast\calO_X(C)) + p_a + 1\\
        &= \deg_C(f^\ast \calO_X(C)) - g_C + 1 - \bigl( \deg_E(g^\ast\calO_X(C)) - g_E + 1\bigr) + p_a + 1\\
        &= C.C - C.E - g_C + g_E + p_a + 1
    \end{align}By Proposition \ref{prop:Adjunction Formula}, we have
    $g_C = \dfrac{1}{2}C.(C+K) + 1, \;g_E = \dfrac{1}{2}E.(E+K) + 1$.
    Thus, substituting them into the above equation and simplifying it, we deduce 
    \begin{align}
        \chi(\calO_X(D)) = \dfrac{1}{2}D.(D-K) + p_a + 1
    \end{align}\hfill
\end{myproof}

\newpage
aaa
\begin{myproof}
    aaaaaaaaaa
\end{myproof}ajfieopfjioapefjeipi\\fjioepaifjeopa
\begin{definition}{}{}
    hogehoge
\end{definition}



\subsection*{Hodge Index Theorem}

\subsection*{Nakai-Moishezon Criterion}
\end{document}